\section{The both-new shell at physical frequency zero}
\label{sec:tpc26-both-new-zero}

The nonzero conductor argument has no information at additive
frequency zero.  We therefore isolate that frequency before introducing
characters or separating the physical row--orbit weight.

\subsection{Uniform endpoint calibration}

For \(M,Y\ge1\), set

\begin{equation}
 A_M(Y)=
 \sum_{\substack{w\le Y\\(w,M)=1}}\frac{a(w)}w,
 \qquad
 B_M(Y)=
 \sum_{\substack{w\le Y\\(w,M)=1}}\frac{\mu(w)}w.
 \label{eq:tpc26-AB}
\end{equation}

We use the following uniform form proved in TPC-24 from the classical
zero-free region for \(\zeta(s)\) \citep{WangTPC24,IwaniecKowalski2004}.

\begin{lemma}[Uniform shell calibration]
\label{lem:tpc26-uniform-calibration}
For every fixed \(A,C>0\), uniformly for \(Y\ge3\) and \(M\le Y^C\),

\begin{align}
 A_M(Y)
 &=\frac{M}{\varphi(M)}
 +O_{A,C}\left(
 \frac{M}{\varphi(M)}(\log Y)^{-A}\right),
 \label{eq:tpc26-A-calibration}\\
 B_M(Y)
 &\ll_{A,C}
 \frac{M}{\varphi(M)}(\log Y)^{-A}.
 \label{eq:tpc26-B-calibration}
\end{align}

Consequently, if \(Y_2\ge Y_1\ge X^\kappa\) and
\(M\le X^{C_0}\), then both endpoint differences

\begin{equation}
 A_M(Y_2)-A_M(Y_1),
 \qquad B_M(Y_2)-B_M(Y_1)
 \label{eq:tpc26-endpoint-differences}
\end{equation}

are \(O_{A,C_0,\kappa}((\log X)^{-A})\) after increasing the input
calibration exponent by a fixed amount.
\end{lemma}

The first main constant \(M/\varphi(M)\) is independent of \(Y\), so
it cancels only after the complete interval is retained.  This is not a
pointwise fixed-modulus statement.

\subsection{Exact gcd reparametrization}

For an off-diagonal row pair \(m\ne n\), put

\begin{equation}
 \Delta=m-n\ne0.
 \label{eq:tpc26-Delta}
\end{equation}

The raw Poisson zero frequency of the both-new joint shell is governed
by

\begin{equation}
 K^{\square,0}_{S,T}(m,n)
 =
 \sum_{\substack{S<u,v\le T\\
                  (u,mH)=(v,nH)=1}}
 \frac{a(u)a(v)}{[u,v]}
 \one_{(u,v)\mid\Delta}.
 \label{eq:tpc26-square-zero-kernel}
\end{equation}

All effective variables are squarefree.

\begin{lemma}[Both-new gcd normal form]
\label{lem:tpc26-both-new-gcd}
The kernel in \eqref{eq:tpc26-square-zero-kernel} equals

\begin{align}
 K^{\square,0}_{S,T}(m,n)
 ={}&
 \sum_{\substack{g\mid\Delta\\g\ \mathrm{squarefree}\\
                  (g,mnH)=1}}
 \frac1g
 \sum_{\substack{S/g<c\le T/g\\(c,gmH)=1}}
 \frac{a(gc)}c
 \notag\\
 &\quad\times
 \sum_{\substack{S/g<d\le T/g\\(d,gcnH)=1}}
 \frac{a(gd)}d.
 \label{eq:tpc26-both-new-gcd}
\end{align}
\end{lemma}

\begin{proof}
For one effective pair write

\[
 g=(u,v),\qquad u=gc,\qquad v=gd.
\]

Squarefree support and exactness of the gcd make \(g,c,d\) pairwise
coprime.  Moreover,

\[
 [u,v]=g\,c\,d,
 \qquad
 \frac1{[u,v]}=\frac1{g\,c\,d}.
\]

Compatibility is precisely \(g\mid\Delta\).  The primitive conditions
become \((g,mnH)=1\), \((c,gmH)=1\), and
\((d,gcnH)=1\).  Substitution gives the formula.
\end{proof}

\subsection{Pointwise cancellation}

\begin{theorem}[Both-new pointwise zero theorem]
\label{thm:tpc26-both-new-zero}
Suppose

\begin{equation}
 S=X^{s+o(1)},
 \qquad T\le X^{C_T},
 \qquad s,C_T>0.
 \label{eq:tpc26-zero-scales}
\end{equation}

For every fixed \(A>0\) and \(0<\kappa<s\), uniformly over the
retained off-diagonal row pairs,

\begin{equation}
 \boxed{
 K^{\square,0}_{S,T}(m,n)
 \ll_{A,\kappa,C_T,h}
 (\log X)^{-A}+X^{-s+\kappa+o(1)}.}
 \label{eq:tpc26-both-new-zero-bound}
\end{equation}
\end{theorem}

\begin{proof}
In \eqref{eq:tpc26-both-new-gcd} split the shared factor at

\begin{equation}
 G_0=SX^{-\kappa}.
 \label{eq:tpc26-G0}
\end{equation}

Suppose first that \(g\le G_0\).  Then \(S/g\ge X^\kappa\).  Fix
\(g,c\), put

\begin{equation}
 M=gcnH,
 \label{eq:tpc26-small-g-M}
\end{equation}

and note that \(M\le X^{C_0}\) for a fixed \(C_0\), uniformly on the
physical fixed-power ranges.  On the effective coprime support,

\begin{equation}
 a(gd)=\mu(g)a(d)-\mu(g)(\log g)\mu(d).
 \label{eq:tpc26-agd-split}
\end{equation}

Hence the innermost sum in \eqref{eq:tpc26-both-new-gcd} is exactly

\begin{align}
 &\mu(g)\{A_M(T/g)-A_M(S/g)\}
 \notag\\
 &\qquad
 -\mu(g)(\log g)\{B_M(T/g)-B_M(S/g)\}.
 \label{eq:tpc26-inner-endpoint}
\end{align}

By \cref{lem:tpc26-uniform-calibration}, this has arbitrary
logarithmic decay.  The outer cofactor mass obeys

\begin{equation}
 \sum_{c\le T/g}\frac{|a(gc)|}{c}
 \ll_{C_T}(\log(2X))^2,
 \label{eq:tpc26-outer-mass}
\end{equation}

and

\[
 \sum_{g\mid\Delta}\frac1g\ll\log\log(3|\Delta|).
\]

Increasing the calibration exponent proves

\begin{equation}
 K^{\square,0,\,g\le G_0}_{S,T}(m,n)
 \ll_A(\log X)^{-A}.
 \label{eq:tpc26-small-g-bound}
\end{equation}

For \(g>G_0\), take absolute values in both cofactor sums.  Each is
bounded by a fixed power of \(\log X\), and the nonzero determinant
gives

\begin{equation}
 \sum_{\substack{g\mid\Delta\\g>G_0}}\frac1g
 \le\frac{\tau(|\Delta|)}{G_0}
 =X^{-s+\kappa+o(1)}.
 \label{eq:tpc26-large-g-bound}
\end{equation}

Combining \eqref{eq:tpc26-small-g-bound} and
\eqref{eq:tpc26-large-g-bound} proves the theorem.
\end{proof}

\begin{remark}[Why the principal row mode disappears here]
\label{rem:tpc26-principal-mode-cancelled}
TPC-25 exhibits a finite shell with a nonzero constant \(g=1\) row
kernel.  In the complete both-new interval, the \(g=1\) coefficient is
an endpoint difference of the form
\eqref{eq:tpc26-inner-endpoint}; its main constant cancels.  Thus the
physical proof does not require projection away from the constant row
vector.  This conclusion depends on retaining the complete shell, not
on a fixed-modulus calibration assertion.
\end{remark}

\subsection{The complete raw annulus}

Define the two mixed kernels

\begin{align}
 K^{ba,0}_{S,T}(m,n)
 :={}&
 \sum_{\substack{u\le S,\ S<v\le T\\
                  (u,mH)=(v,nH)=1}}
 \frac{b_R(u)a(v)}{[u,v]}
 \one_{(u,v)\mid\Delta},
 \label{eq:tpc26-ba-kernel}\\
 K^{ab,0}_{S,T}(m,n)
 :={}&K^{ba,0}_{S,T}(n,m).
 \label{eq:tpc26-ab-kernel}
\end{align}

The raw annular joint zero kernel is

\begin{equation}
 K^{\rm ann,0}_{S,T}(m,n)
 =K^{ba,0}_{S,T}(m,n)+K^{ab,0}_{S,T}(m,n)
 +K^{\square,0}_{S,T}(m,n).
 \label{eq:tpc26-annular-zero-decomposition}
\end{equation}

TPC-24 proves the analogue of
\eqref{eq:tpc26-both-new-zero-bound} for \(K^{ba,0}\), and transposition
gives it for \(K^{ab,0}\).  We therefore obtain:

\begin{corollary}[Annular raw zero kernel]
\label{cor:tpc26-annular-raw-zero}
Under the hypotheses of \cref{thm:tpc26-both-new-zero},

\begin{equation}
 \boxed{
 K^{\rm ann,0}_{S,T}(m,n)
 \ll_{A,\kappa,C_T,h}
 (\log X)^{-A}+X^{-s+\kappa+o(1)}.}
 \label{eq:tpc26-annular-raw-zero-bound}
\end{equation}
\end{corollary}

The two single-density drift terms from
\cref{prop:tpc26-annular-identity} are not part of
\eqref{eq:tpc26-annular-zero-decomposition}; they remain explicit and
are controlled by \cref{lem:tpc26-drift-bound}.
